% ===================================================================
%  TAULUKKO N:o 6 — Talo sisältä, huonekalut, ravinto, valaistus.
%  Traduction 2026 d'apres texto/io, controlee sur texto/fr et
%  texto/sv. Consigne : voir texto/fi/10-taulukko-01.tex.
%
%  Deux notes, toutes deux marquees « (*) », et deux appels.
%  LA FEMME DE CHAMBRE PREND SON (16), contre le fac-simile ido :
%  ecart declare dans outils/renvoji.py.
% ===================================================================

%%K t06-apar-1 apar
\VUsaut{6.0mm}
\VUtitre{15.0pt}{60}{TAULUKKO N:o 6}
\VUsaut{1.4mm}
\VUtitre{11.4pt}{0}{Talo sisältä, huonekalut, ravinto,}
\VUsaut{0.4mm}
\VUtitre{11.4pt}{0}{valaistus.}
\VUsaut{2.2mm}

%%K t06-apar-2 apar
Talo on rakennettu valmiiksi. Juhanin setä on asettunut uuteen asuntoonsa, jonka
jaon tunnemme. Katsokaamme nyt, kuinka hän on kalustanut talonsa. Tutkimme
ensin:

%%K t06-tit-1 sub
\VUpk{10.2pt}{40}{Makuuhuone.}
\VUsaut{3.0mm}

%%K t06-c1-01-1 p
1. --- Tulemme sopimattomaan aikaan, sillä kamarineitsyt on siivoamassa
\VUgras{huonetta}.

%%K t06-c1-02-1 p
2. --- \VUgras{Pesupöydällä} \textsuperscript{(1)} on siellä täällä eri
esineitä, joilla peseydytään, kammataan, sanalla sanoen hoidetaan ulkonäköä. Ne
ovat \VUgras{pesuvati} \textsuperscript{(2)} ja \VUgras{vesikannu}
\textsuperscript{(3)}, \VUgras{hiusharja} \textsuperscript{(4)},
\VUgras{kampa} \textsuperscript{(5)}, \VUgras{saippua} \textsuperscript{(6)} ja
\VUgras{sieni} \textsuperscript{(7)}, lopuksi pieni \VUgras{pullo}
\textsuperscript{(8)} nestemäistä hammasainetta varten.

%%K t06-c1-03-1 p
3. --- Lattialla, pesupöydän edessä, tässä on \VUgras{sanko}
\textsuperscript{(9)} likaisen veden keräämistä varten.

%%K t06-c1-04-1 p
4. --- Näemme \VUgras{eteisessä} \textsuperscript{(10)} muutamia
\VUgras{naulakoita} \textsuperscript{(13)} ja kaksi \VUgras{sateenvarjoa}
\textsuperscript{(11)} \VUgras{sateenvarjotelineessä} \textsuperscript{(12)}.

%%K t06-c1-05-1 p
5. --- Oven toisella puolella on \VUgras{peilikaappi} \textsuperscript{(14)},
joka sisältää \VUgras{liinavaatteet} \textsuperscript{(15)}.

%%K t06-c1-06-1 p
6. --- Käyttäkäämme hyväksemme sitä, että \VUgras{kamarineitsyt}
\textsuperscript{(16)} sijaa \VUgras{sänkyä} \textsuperscript{(17)},
tutkiaksemme sen eri osat. \VUgras{Sängynrunko} \textsuperscript{(20)} sisältää
hyvän \VUgras{joustinpatjan} \textsuperscript{(21)}, jolle Justiina heti panee
villaisen \VUgras{patjan} \textsuperscript{(22)}. Sitten hän ottaa tuoleilta
kaksi \VUgras{lakanaa} \textsuperscript{(23)} ja \VUgras{huovan}
\textsuperscript{(24)}. Silloin hän panee päänpuoleen \VUgras{poikittaisen
tyynyn} \textsuperscript{(25)} ja \VUgras{päätyynyn} \textsuperscript{(26)},
jotka ovat \VUgras{untuvapeitteen} \textsuperscript{(27)} kanssa
\VUgras{vuodematolla} \textsuperscript{(32)}. Hän käyttää höyhenhuiskaa
pölyttääkseen \VUgras{verhot} \textsuperscript{(19)} ja \VUgras{sängyn katoksen}
\textsuperscript{(18)}, ja tässä on osa työstä valmis.

%%K t06-c1-07-1 p
7. --- Silloin hänen täytyy järjestää \VUgras{yöpöytää} \textsuperscript{(28)}
hieman, vetää \VUgras{herätyskello} \textsuperscript{(29)}, valmistaa
\VUgras{yölamppu} \textsuperscript{(30)}, poistaa \VUgras{kynttilä}
\textsuperscript{(31)} puhdistaakseen kynttilänjalan.

%%K t06-tit-2 sub
\VUsaut{5.0mm}
\VUpk{10.2pt}{40}{Ompelukori ja takan varusteet.}
\VUsaut{3.0mm}

%%K t06-c1-08-1 p
8. --- \VUgras{Ompelukori} \textsuperscript{(34)} \VUgras{jakkaran}
\textsuperscript{(33)} vieressä tarvitsee myös kovasti järjestämistä. Hänen
täytyy taittaa kokoon \VUgras{kirjailut} \textsuperscript{(35)}, sulkea
\VUgras{ompelupöytään} \textsuperscript{(38)} \VUgras{sormustin},
\VUgras{lanka} \textsuperscript{(36)}, \VUgras{neulat} \textsuperscript{(37)} ja
\VUgras{sakset} \textsuperscript{(39)} ja sulkea pöydän kansi
\VUgras{avaimella} \textsuperscript{(40)}.

%%K t06-c1-09-1 p
9. --- Keskellä olevalla \VUgras{yksijalkaisella pöydällä}
\textsuperscript{(41)} Justiina vaihtaa veden \VUgras{karahvissa}
\textsuperscript{(42)}. Hän panee \VUgras{sokeria}
\VUgras{sokerikkoon} \textsuperscript{(43)}, puhdistaa \VUgras{sokeripihdit}
\textsuperscript{(44)} ja poistaa \VUgras{vaateharjan} \textsuperscript{(46)}.

%%K t06-c1-10-1 p
10. --- Huoneen oikea puoli vaatii häneltä vähemmän työtä, sillä
\VUgras{lepotuoli} \textsuperscript{(47)} ja sen \VUgras{tyyny}
\textsuperscript{(48)} ovat oikealla paikallaan, ja koska sää ei ole kylmä,
mitään ei ole siirretty \VUgras{takan} \textsuperscript{(49)} varusteiden
(välineiden) joukossa. \VUgras{Lapio} \textsuperscript{(50)},
\VUgras{pihdit} \textsuperscript{(51)}, \VUgras{tulipukit}
\textsuperscript{(55)}, \VUgras{tuhkaritilä} \textsuperscript{(56)} ovat
puhtaat; \VUgras{tulisuojusta} \textsuperscript{(54)} ei ole siirretty ja
\VUgras{halkolaatikko} \textsuperscript{(52)} on hyvin varustettu
\VUgras{polttopuilla} \textsuperscript{(53)}.

%%K t06-noto-1 noto
\VUnotes{143.2mm}{%
(*) Babo = pieni lapsi, engl. baby, ransk. bébé, ital. bambino.%
}

%%K t06-noto-2 noto
\VUnotes{143.2mm}{%
(*) Lambrequino = ransk. lambrequin, esp. lambrequines, port. lambrequins, jopa
saks. engl. lambrequins, siis saks. engl. ransk. port. esp.%
}

%%K t06-c1-11-1 p
11. --- Kuitenkin hänen täytyy pölyttää \VUgras{kaappikello}
\textsuperscript{(57)}, \VUgras{kynttiläkruunut} \textsuperscript{(58)} ja
\VUgras{pikku lippaat} \textsuperscript{(59)}, lopuksi pyyhkiä \VUgras{peili}
\textsuperscript{(60)}.

%%K t06-c1-12-1 p
12. --- Mitä pienen lapsen (vauvan (*)) \VUgras{kehtoon}
\textsuperscript{(61)} tulee, se on jo valmiina.

%%K t06-tit-3 sub
\VUsaut{5.0mm}
\VUpk{10.2pt}{40}{Sali.}
\VUsaut{3.0mm}

%%K t06-c2-01-1 p
1. --- Nyt tässä on sali. Se on hyvin kaunis \VUgras{huone}. Takka, joka on
oikeassa nurkassa, on koristettu hienolla \VUgras{pienoispatsaalla}
\textsuperscript{(1)}, kahdella kauniilla \VUgras{maljakolla}
\textsuperscript{(3)} ja verhoiltu samettisella \VUgras{reunusverholla}
\textsuperscript{(2)} (*).

%%K t06-c2-02-1 p
2. --- Estääkseen sitä, että tulisijasta lentävät kipinät sytyttäisivät maton, on
tulisijan eteen asetettu kuparinen \VUgras{kipinäritilä} \textsuperscript{(4)}.

%%K t06-c2-03-1 p
3. --- Aivan lähellä on tyylikäs naisten \VUgras{kirjoituspöytä}
\textsuperscript{(5)} avoinna. Sen ääressä \VUgras{talon rouva}
\textsuperscript{(23)} kirjoittaa kirjeitään.

%%K t06-c2-04-1 p
4. --- Ikkuna on varustettu \VUgras{ohuilla verhoilla} \textsuperscript{(6)},
joissa on \VUgras{nauhat} \textsuperscript{(8)} kiinnitettyinä
\VUgras{koukkuihin} \textsuperscript{(7)}, ja kangasverhoilla, jotka on
\VUgras{ylhäältä laskostettu} \textsuperscript{(9)}.

%%K t06-c2-05-1 p
5. --- Ikkunasyvennyksessä \VUgras{ruukunsuojus} \textsuperscript{(11)} sisältää
vihreän \VUgras{kasvin} \textsuperscript{(10)}, komean palmun, jonka lehdet
ulottuvat miltei \VUgras{öljylamppuun} \textsuperscript{(12)} asti.

%%K t06-c2-06-1 p
6. --- Lampun \VUgras{varjostimen} \textsuperscript{(13)} on järjestänyt Maria,
se hurmaava \VUgras{nuori tyttö} \textsuperscript{(14)}, joka istuu pianon
\textsuperscript{(15)} ääressä. Istuen hyvin suorassa \VUgras{jakkaralla}
\textsuperscript{(16)} hän harjoittelee musiikkikappaletta. Hän on hyvin taitava
ja huolellinen, kuten hänen \VUgras{nuottikaappinsa} \textsuperscript{(17)}
todistaa, joka on täydellisesti järjestetty.

%%K t06-tit-4 sub
\VUsaut{5.0mm}
\VUpk{10.2pt}{40}{Vintiö.}
\VUsaut{3.0mm}

%%K t06-c2-07-1 p
7. --- Hänen veljensä Paavali etsii \VUgras{kirjaa} \textsuperscript{(19)}
\VUgras{kirjahyllystä} \textsuperscript{(18)}. Hän on \VUgras{nuori poika}
\textsuperscript{(20)}, vilkas ja sietämätön. Kun hän riippuu kiinni
kirjahyllyssä yltääkseen kirjaan, joka on liian korkealla, hän on vaarassa
pudottaa \VUgras{rintakuvan} \textsuperscript{(21)}, joka on sen päällä.

%%K t06-c2-08-1 p
8. --- Hän käyttää hyväkseen sitä, että hänen äitinsä on kiireinen
keskustellessaan ystävättären kanssa, \VUgras{rouvan} \textsuperscript{(22)},
joka tuli viettämään muutaman päivän hänen luonaan, tehdäkseen tuon
tyhmyyden. Äiti istuu \VUgras{sohvalla} \textsuperscript{(24)}
\VUgras{nojatuolin} \textsuperscript{(25)} vieressä, ja hänen ystävättärensä
istuu \VUgras{tuolilla} \textsuperscript{(27)}, josta näkyy vain
\VUgras{selkänoja} \textsuperscript{(26)}.

%%K t06-c2-09-1 p
9. --- Suuri \VUgras{kehys} \textsuperscript{(30)}, joka riippuu seinällä,
sisältää erään sukulaisen \VUgras{muotokuvan} \textsuperscript{(29)}.
\VUgras{Albumiin} \textsuperscript{(32)}, joka on pöydällä, on koottu muiden
perheenjäsenten \VUgras{valokuvat} \textsuperscript{(33)}. Lapset ovat juuri
selailleet sitä, sillä \VUgras{tuolit} \textsuperscript{(31)} ovat vielä
ryhmiteltyinä pöydän ympärille.

%%K t06-c2-10-1 p
10. --- \VUgras{Kiinalainen maljakko} \textsuperscript{(34)}, joka on posliinia
ja joka on \VUgras{paperiveitsen} \textsuperscript{(35)} vieressä, lahjoitettiin
Marialle hänen nimipäiväkseen, ja \VUgras{taittoseinän} \textsuperscript{(36)},
joka koristaa huoneen nurkkaa, toi Kiinasta hänen setänsä.

%%K t06-c2-11-1 p
11. --- Elävöitettynä komeilla \VUgras{maalauksilla} \textsuperscript{(37)},
jotka riippuvat \VUgras{seinällä} \textsuperscript{(42)}, valaistuna
\VUgras{kattokruunulla} \textsuperscript{(38)}, jonka \VUgras{sähkölamput}
\textsuperscript{(39)} loistavat iloisesti illalla, tuo sali näyttää hyvin
miellyttävältä. Pehmeä \VUgras{matto} \textsuperscript{(40)}, joka on levitetty
parketille, ja niin helppokäyttöinen \VUgras{sähkösoittokello}
\textsuperscript{(41)} täydentävät sen mukavuuden.

%%K t06-tit-5 sub
\VUsaut{3.6mm}
\VUpk{10.2pt}{40}{Ruokasali.}
\VUsaut{3.0mm}

%%K t06-c3-01-1 p
1. --- Päivälliskello on soinut. Koko perhe, Mariaa lukuun ottamatta, on koolla
pöydän ympärillä. Ruokasalia lämmittää miellyttävästi erinomainen fajanssinen
\VUgras{kamiina} \textsuperscript{(1)}, jossa on ohut \VUgras{putki}
\textsuperscript{(2)}.

%%K t06-c3-02-1 p
2. --- Tuossa huoneessa ei ole monta huonekalua. Seinustalla riippuu,
\VUgras{penkin} \textsuperscript{(4)} yläpuolella, sievä \VUgras{vesisäiliö}
\textsuperscript{(3)}. Aivan vieressä on \VUgras{senkki}
\textsuperscript{(5)}, jolle asetetaan kylmät ruoat, jälkiruokavadit ja ne eri
pöytävälineet, joita ei tarvita heti.

%%K t06-c3-03-1 p
3. --- Niinpä näemme ylähyllyllä \VUgras{paistetun kanan} \textsuperscript{(7)},
\VUgras{kalan} \textsuperscript{(8)} ja \VUgras{maljan} \textsuperscript{(10)},
jossa on \VUgras{hedelmiä} \textsuperscript{(9)}. Alempana tässä ovat
\VUgras{juusto} \textsuperscript{(11)}, \VUgras{leipä} \textsuperscript{(12)},
\VUgras{pöytäteline} \textsuperscript{(13)}, joka sisältää kaiken mitä tarvitaan
salaatin maustamiseen: öljyn, etikan, \VUgras{suolan} \textsuperscript{(14)} ja
\VUgras{pippurin} \textsuperscript{(15)}. Lopuksi alahyllyllä on \VUgras{kakku}
\textsuperscript{(17)} \VUgras{sinappipurkin} \textsuperscript{(16)} vieressä.

%%K t06-c3-04-1 p
4. --- Ruokasalin kellolla, jota sanotaan \VUgras{seinäkelloksi}
\textsuperscript{(20)}, on hyvin omalaatuinen muoto. Se on upotettu
puukehykseen, joka sisältää lisäksi \VUgras{ilmapuntarin}
\textsuperscript{(19)} ja \VUgras{lämpömittarin} \textsuperscript{(18)}.

%%K t06-tit-6 sub
\VUsaut{4.6mm}
\VUpk{10.2pt}{40}{Päivällisen tarjoilu.}
\VUsaut{3.0mm}

%%K t06-c3-05-1 p
5. --- Huoneen kaikilla seinillä on \VUgras{puupaneeleja}
\textsuperscript{(21)} vahattua tammea. Kankainen \VUgras{oviverho}
\textsuperscript{(22)} sulkee riittävästi oven, joka johtaa kuistille. Sinne
perhe menee juomaan kahvin päivällisen jälkeen. Jo ovat \VUgras{kupit}
\textsuperscript{(24)} ja \VUgras{lautaset} \textsuperscript{(25)} valmiina
\VUgras{tarjoilupöydällä} \textsuperscript{(23)}.

%%K t06-c3-06-1 p
6. --- Pöydän yläpuolelle on \VUgras{riippulamppu} \textsuperscript{(26)}
kiinnitetty kattoon, ja sen hyvin kirkas lamppu valaisee illalla koko
ruokasalin.

%%K t06-c3-07-1 p
7. --- Tutkikaamme nyt itse \VUgras{ruokapöytää} \textsuperscript{(27)}. Kun
perhe tuli sisään, kattaus oli valmis. \VUgras{Pöytäliinalle}
\textsuperscript{(28)} oli kunkin pöytävieraan paikalle asetettu
\VUgras{lautanen} \textsuperscript{(33)}, \VUgras{lautasliina}
\textsuperscript{(29)}, \VUgras{lusikka} \textsuperscript{(34)},
\VUgras{haarukka} \textsuperscript{(35)}, \VUgras{veitsi}
\textsuperscript{(36)} ja \VUgras{lasi} \textsuperscript{(32)}. Oli myös
\VUgras{pulloja} \textsuperscript{(30)} viiniä varten ja \VUgras{karahvi}
\textsuperscript{(31)} vettä varten.

%%K t06-c3-08-1 p
8. --- \VUgras{Palvelija} \textsuperscript{(39)} kantoi ensin sisään
\VUgras{liemimaljan} \textsuperscript{(37)}, jossa vielä höyryää hieman keittoa.
Nyt hän tarjoilee ensimmäisen \VUgras{ruokalajin} \textsuperscript{(38)},
liharuoan. Sitten hän tarjoaa niitä, jotka olemme jo maininneet, lopuksi,
\VUgras{jälkiruokavadin} jälkeen, hän käyttää hyväkseen sitä, että hänen
isäntäväkensä on mennyt kuistille, poistaakseen pöytävälineet ja järjestääkseen
ruokasalin uudelleen.

%%K t06-c3-08-2 p
\textit{Huom.} --- \textit{Ravintoa koskevat yleiset sanat annetaan tässä hyvin
lyhyesti. Luettelo täydennetään kahdessa taulukossa « Hotelli » (n:o 13) ja
« Markkinat » (n:o 15).}

%%K t06-tit-7 sub
\VUsaut{4.6mm}
\VUpk{10.2pt}{40}{Keittiö.}
\VUsaut{3.0mm}

%%K t06-c4-01-1 p
1. --- Keittiössä, jota alamme katsella puoliavoimesta ovesta, näemme
maalauksellista epäjärjestystä. \VUgras{Tulisijan} \textsuperscript{(1)} edessä,
jossa \VUgras{tuli} \textsuperscript{(3)} rätisee, on \VUgras{paistinvarras}
\textsuperscript{(5)}. Kaunis \VUgras{paisti} \textsuperscript{(7)} on
\VUgras{vartaaseen pistettynä} \textsuperscript{(6)} ja tulee valmiiksi.

%%K t06-c4-02-1 p
2. --- \VUgras{Palkeet} \textsuperscript{(8)} ja \VUgras{hiililapio}
\textsuperscript{(9)}, joita on juuri käytetty, ovat jääneet lattialle samoin
kuin \VUgras{polttopuut} \textsuperscript{(10)}.

%%K t06-c4-03-1 p
3. --- Takanreunus on järjestetty; \VUgras{lyhty} \textsuperscript{(11)},
\VUgras{tulitikkurasia} \textsuperscript{(13)}, jossa \VUgras{tulitikut}
\textsuperscript{(12)} ovat, \VUgras{kynttilänjalka} \textsuperscript{(14)} ja
\VUgras{käsikynttilänjalka} \textsuperscript{(15)}, jossa on \VUgras{kynttilä}
\textsuperscript{(16)}, ovat paikoillaan. Mutta suuri \VUgras{hiilisanko}
\textsuperscript{(18)}, täynnä \VUgras{kivihiiltä} \textsuperscript{(17)}, jonka
\VUgras{keittäjätär} \textsuperscript{(23)} on juuri täyttänyt kellarissa, on
jäänyt keskelle keittiötä.

%%K t06-c4-04-1 p
4. --- Totisesti, poloisen naisen täytyy kiirehtiä, jotta aamiainen olisi valmis
oikeaan aikaan. Nyt hän on \VUgras{hellan} \textsuperscript{(19)} ääressä ja
hämmentää \VUgras{kastiketta} \textsuperscript{(24)}, joka ei saa palaa pohjaan
liikaa.

%%K t06-c4-05-1 p
5. --- \VUgras{Uunissa} \textsuperscript{(20)} paistuu kakku, jonka hän on
valmistanut lasten välipalaksi. Hellan yläpuolella riippuvat \VUgras{kauha}
\textsuperscript{(21)} ja \VUgras{reikäkauha} \textsuperscript{(22)}.

%%K t06-tit-8 sub
\VUsaut{4.0mm}
\VUpk{10.2pt}{40}{Talousastiat.}
\VUsaut{3.0mm}

%%K t06-c4-06-1 p
6. --- \VUgras{Keittoastiat} \textsuperscript{(25)}, jotka riippuvat huoneen
perällä, ovat hyvin puhtaat. Kauniit kupariset \VUgras{kattilat}
\textsuperscript{(26)}, \VUgras{maitokannu} \textsuperscript{(27)}, pieni
\VUgras{siivilä} \textsuperscript{(28)} ja \VUgras{raastinrauta}
\textsuperscript{(29)} on huolellisesti puhdistettu.

%%K t06-c4-07-1 p
7. --- \VUgras{Suola-astia} \textsuperscript{(30)} on senkillä
\VUgras{teekannun} \textsuperscript{(31)} ja \VUgras{suuren öljypullon}
\textsuperscript{(32)} vieressä. \VUgras{Laatikko} \textsuperscript{(33)}
sisältää luultavasti ruokailuvälineet ja keittiön veitset.

%%K t06-c4-08-1 p
8. --- \VUgras{Luuta} \textsuperscript{(34)} ja \VUgras{höyhenhuiska}
\textsuperscript{(35)} ovat nurkassa. Keittäjätär on juuri käyttänyt
\VUgras{hakkuupölkkyä} \textsuperscript{(36)}, hän on jättänyt sen ikkunan
viereen.

%%K t06-c4-09-1 p
9. --- \VUgras{Pyyheliina} \textsuperscript{(37)} riippuu muurilla,
\VUgras{suppilon} \textsuperscript{(38)} vieressä. Keittiön siinä osassa ovat
\VUgras{halstari} \textsuperscript{(39)}, \VUgras{hakkuuveitsi}
\textsuperscript{(40)} ja \VUgras{hakkuulauta} \textsuperscript{(41)}
sijoitettuina. Sitten ovat, hakkuuveitsen yläpuolella, \VUgras{hyllyllä}
\textsuperscript{(42)}, \VUgras{ruukkuja} \textsuperscript{(43)},
\VUgras{keittopata} \textsuperscript{(44)} \VUgras{kansineen}
\textsuperscript{(45)} ja \VUgras{kattila} \textsuperscript{(46)}.
\VUgras{Paistinpannu} \textsuperscript{(47)} riippuu aivan lähellä.

%%K t06-c4-10-1 p
10. --- Vain kuparinen \VUgras{hana} \textsuperscript{(48)} luo kiiltoa tuohon
nurkkaan. \VUgras{Tiskiallas} \textsuperscript{(49)}, jolla paksu \VUgras{suuri
kannu} \textsuperscript{(50)} seisoo, on luultavasti hyvin mukava pienine
kaappeineen, jossa säilytetään \VUgras{etikka-astiaa} \textsuperscript{(51)},
joka on varustettu \VUgras{tapillaan} \textsuperscript{(52)},
\VUgras{jätelaatikkoa} \textsuperscript{(53)} ja \VUgras{likaisia lautasia}
\textsuperscript{(54)}.

%%K t06-tit-9 sub
\VUsaut{4.0mm}
\VUpk{10.2pt}{40}{Muita esineitä keittiössä.}
\VUsaut{3.0mm}

%%K t06-c4-11-1 p
11. --- \VUgras{Saavi} \textsuperscript{(55)}, jossa pestään
\VUgras{vihannekset} \textsuperscript{(58)}, ja valurautainen \VUgras{pata}
\textsuperscript{(56)}, joka on \VUgras{kivilattialla} \textsuperscript{(57)},
ovat luultavasti myös sijoitetut tuohon kaappiin.

%%K t06-c4-12-1 p
12. --- Keittäjättäreltä ei kaiketi puutu liesiä, sillä tässä on lisäksi, pöydän
nurkalla, \VUgras{kaasukeitin} \textsuperscript{(59)}, jolla lämmitetään vettä
pienessä kattilassa. \VUgras{Höyry} \textsuperscript{(60)}, joka pääsee siitä
ulos, osoittaa meille, että se luultavasti kiehuu. Todennäköisesti se vesi
käytetään kahviin, sillä tässä on pöydän toisessa päässä
\VUgras{kahvipannu} \textsuperscript{(64)} ja \VUgras{kahvimylly}
\textsuperscript{(63)}.

%%K t06-c4-13-1 p
13. --- Keittiö on yhdistetty \VUgras{kaasupolttimeen} \textsuperscript{(62)}
pitkällä \VUgras{kumiletkulla} \textsuperscript{(61)}.

%%K t06-c4-14-1 p
14. --- \VUgras{Päivittäisapulainen} \textsuperscript{(66)}, joka tulee
auttamaan keittäjätärtä, valmistaa vihannekset päivälliseksi. Kun ne on
puhdistettu ja pesty, hän panee ne \VUgras{kulhoon} \textsuperscript{(65)}
odottamaan niiden keittämisen hetkeä.

%%K t06-tit-10 sub
\VUsaut{4.0mm}
{\centering\fontsize{10.2pt}{10.2pt}\selectfont\textls[40]{\scshape Kylpyhuone.} \textit{(Kylpykammari.)}\par}
\VUsaut{3.0mm}

%%K t06-c5-01-1 p
1. --- Meille jää vain yksi hienotunteettomuus tehtäväksi oppiaksemme tuntemaan
asunnon tärkeimmät huoneet: nähdä kylpyhuoneen sisustus. Se ei vaadi pitkää
aikaa, sillä se ei ole suuri, ja saamme nopeasti tietää, että
\VUgras{kylpyammeessa} \textsuperscript{(1)} on hyvä \VUgras{vedenlämmitin}
\textsuperscript{(2)}, että \VUgras{saippuakuppi} \textsuperscript{(3)} on
kylpijän käden ulottuvilla ja että \VUgras{pyyheteline}
\textsuperscript{(5)} varmasti antaa \VUgras{pyyhkeiden} \textsuperscript{(4)}
kuivua helposti.

%%K t06-c5-02-1 p
2. --- Tuon huoneen seinät on peitetty \VUgras{kaakeleilla}
\textsuperscript{(6)}, mikä teki mahdolliseksi asentaa \VUgras{suihkulaitteen}
\textsuperscript{(7)} pelkäämättä mitenkään, että vesi, joka roiskuu takaisin
sitä käytettäessä, vahingoittaisi muureja. Pieni sinkkinen \VUgras{vati}
\textsuperscript{(8)}, joka palvelee jalkakylpyihin, täydentää kalustuksen.
\VUsaut{6.0mm}
\VUfilet{22mm}
